\subsection{Upper integrals of positive functions with respect to an integral of point measures}

We are going to see that, when \((\pi,g)\) is a \(\mu\) -adapted pair, one can sharpen the results obtained by applying the propositions of §3 to the family \(t\mapsto\lambda_{t}=g(t)\varepsilon_{\pi(t)}\) , which is \(\mu\) -adequate by Prop. 1.

THEOREM 1. --- Let \((\pi, g)\) be a \(\mu\) -adapted pair, and let

\[
\nu = \int g (t) \varepsilon_ {\pi (t)} d \mu (t).
\]

For every numerical function \(f \geqslant 0\) defined on \(X\) ,

\[
\int^ {\bullet} f (x) d \nu (x) = \int^ {\bullet} f (\pi (t)) g (t) d \mu (t).\tag{1}
\]

\begin{enumerate}
\def\labelenumi{\Alph{enumi})}
\tightlist
\item
  Suppose first that the measure \(\mu\) has compact support \(K\) and that the restrictions to \(K\) of the functions \(g\) and \(\pi\) are continuous. By formula (4) of §3, No.~1, \(\nu^{\bullet}(1) = \int_{K} g(t) d\mu(t) < +\infty\) , so that all of the measures that figure in formula (1) are bounded. We may therefore replace \(\int^{\bullet}\) by \(\int^{*}\) in the first member. In view of formula (6) of §3, No.~2, it all comes down to proving that
\end{enumerate}

\[
\int^ {*} f (x) d \nu (x) \leqslant \int^ {\bullet} f (\pi (t)) g (t) d \mu (t),\tag{2}
\]

where the symbol \(\int^{\bullet}\) in the second member can in turn be replaced by \(\int^{*}\) . By the definition of upper integral, it suffices to verify the inequality

\[
\int^ {*} f (x) d \nu (x) \leqslant \int^ {*} h (t) d \mu (t)\tag{3}
\]

for every lower semi-continuous function \(h\) on \(\mathrm{T}\) that is \(\geqslant\) the function \(t \mapsto f(\pi(t))g(t)\) . Now, let \(\varepsilon\) be a number \(> 0\) and let \(u\) be the function \((h + \varepsilon)/g\) , which is lower semi-continuous in \(\mathrm{K}\) . If \(t \in \overline{\pi}^{-1}(\{x\}) \cap \mathrm{K}\) , then \(u(t) \geqslant f(x)\) : this is obvious if \(g(t) = 0\) , because then \(u(t) = +\infty\) ; if \(g(t) > 0\) , then

\[
u (t) g (t) = h (t) + \varepsilon \geqslant f (\pi (t)) g (t) = f (x) g (t),
\]

whence the asserted inequality. Under these conditions, for every \(x \in X\) let \(v(x)\) be the infimum of \(u(t)\) for \(t \in \overline{\pi}^{-1}(\{x\}) \cap K\) . The function v is \(\geqslant f\) by the foregoing, it is lower semi-continuous on X by the Lemma (applied to the restriction of \(\pi\) to K), and \(v(\pi(t))g(t) \leqslant h(t) + \varepsilon\) for all \(t \in K\) (recall that the first member is zero by convention if \(g(t) = 0\) ). Let us then apply to v the formula (4) of §3, No.~1. We obtain:

\[
\begin{array}{l} \int^ {*} f (x) d \nu (x) \leqslant \int^ {*} v (x) d \nu (x) \\ = \int^ {*} v (\pi (t)) g (t) d \mu (t) \leqslant \int_ {K} ^ {*} (h (t) + \varepsilon) d \mu (t) \\ = \int^ {*} h (t) d \mu (t) + \varepsilon \mu (1). \end{array}\tag{4}
\]

Since the measure \(\mu\) is bounded and \(\varepsilon\) is arbitrary, the inequality (3) follows.

\begin{enumerate}
\def\labelenumi{\Alph{enumi})}
\setcounter{enumi}{1}
\tightlist
\item
  Let us now pass to the general case. Since the mapping \(t \mapsto (\pi(t), g(t))\) of \(T\) into \(X \times R_{+}\) is \(\mu\) -measurable (Ch. IV, §5, No.~3, Th. 1), the set \(\mathfrak{K}\) of compact subsets \(K\) of \(T\) such that the restrictions of \(\pi\) and \(g\) to \(K\) are continuous is \(\mu\) -dense (Ch. IV, §5, No.~10, Prop. 15). By Prop. 4 of §2, No.~3, \(\mu\) is the sum of a summable family \((\mu_{\alpha})_{\alpha \in A}\) of measures whose supports are elements of \(\mathfrak{K}\) ; the pair \((\pi, g)\) being \(\mu_{\alpha}\) -adapted for every \(\alpha \in A\) , let \(\nu_{\alpha}\) be the measure \(\int g(t)\varepsilon_{\pi(t)} d\mu_{\alpha}(t)\) . Then by A),
\end{enumerate}

\[
\int^ {\bullet} f (x) d \nu_ {\alpha} (x) = \int^ {\bullet} f (\pi (t)) g (t) d \mu_ {\alpha} (t).\tag{5}
\]

But the \(\nu_{\alpha}\) form a summable family whose sum is equal to \(\nu\) (§3, No.~1, Cor. of Prop. 1). Therefore, by Prop. 1 of §2, No.~2,

\[
\int^ {\bullet} f (x) d \nu (x) = \sum_ {\alpha \in \mathrm{A}} \int^ {\bullet} f (x) d \nu_ {\alpha} (x).\tag{6}
\]

One has an analogous relation for the second member of (5), and (1) therefore follows from (5) by summing over \(\alpha\) .

COROLLARY. --- In order that a subset N of X be locally negligible for \(\nu\) , it is necessary and sufficient that the intersection of \(\overline{\pi}^{-1}(N)\) with the set of points \(t \in T\) where \(g(t) > 0\) be locally negligible for \(\mu\) .

PROPOSITION 2. --- Let \(\pi\) be a proper continuous mapping (GT, I, §10, No.~1) of T into X, and let \(g\) be a continuous, finite numerical function on T such that \(g(t) > 0\) for all \(t \in \mathrm{T}\) . The pair \((\pi, g)\) is then \(\mu\) -adapted, and if one sets

\[
\nu = \int g (t) \varepsilon_ {\pi (t)} d \mu (t),
\]

then, for every numerical function \(f \geqslant 0\) defined on \(X\) ,

\[
\int^ {*} f (x) d \nu (x) = \int^ {*} f (\pi (t)) g (t) d \mu (t).\tag{7}
\]

It is clear that \(\pi\) and g are \(\mu\) -measurable; moreover, for every function \(\psi \in \mathcal{K}(\mathrm{X})\) , \(\psi \circ \pi\) is continuous with compact support, since \(\pi\) is proper; the pair \((\pi, g)\) is therefore \(\mu\) -adapted and, moreover, the mapping \(t \mapsto g(t) \varepsilon_{\pi(t)}\) is vaguely continuous.

Let h be a lower semi-continuous function on T such that

\[
f (\pi (t)) g (t) \leqslant h (t) \quad \text { for   all } t \in \mathrm{T}.
\]

We are going to show that

\[
\int^ {*} f (x) d \nu (x) \leqslant \int^ {*} h (t) d \mu (t).\tag{8}
\]

By the definition of upper integral, this will imply the inequality

\[
\int^ {*} f (x) d \nu (x) \leqslant \int^ {*} f (\pi (t)) g (t) d \mu (t),
\]

which, combined with the inequality (7) of §3, No.~2, will prove (7).

To prove (8), let us define a function \(\overline{f}\) on \(X\) in the following manner: \(\overline{f}(x)\) is the infimum of \(h(t) / g(t)\) in the set \(\overline{\pi}^{-1}(x)\) (infimum equal to \(+\infty\) if \(\overline{\pi}^{-1}(x) = \varnothing\) ). The function \(\overline{f}\) has the following properties:

\(1^{\circ}\overline{f} (x)\geqslant f(x)\) for all \(x\in \mathrm{X}\) (since \(g(t) > 0\) for all \(t\in \mathrm{T}\) ).

\[
2 ^ {\circ} \overline {{f}} (\pi (t)) g (t) \leqslant h (t) \text {   for   all   } t \in \mathrm{T}.
\]

\(3^{\circ}\) The function \(\overline{f}\) is lower semi-continuous by virtue of the Lemma, the function \(h / g\) being lower semi-continuous on T.

Consequently, in view of Prop. 2a) of §3, No.~1:

\[
\int^ {*} f (x) d \nu (x) \leqslant \int^ {*} \overline {{{f}}} (x) d \nu (x) = \int^ {*} \overline {{{f}}} (\pi (t)) g (t) d \mu (t) \leqslant \int^ {*} h (t) d \mu (t),
\]

which establishes (8), and completes the proof.

